\documentclass[12pt]{article}
\usepackage[utf8]{inputenc}
\usepackage[T1]{fontenc}
\usepackage[margin=1in]{geometry}
\usepackage{setspace}
\usepackage{xcolor}
\usepackage[colorlinks=true, urlcolor=blue!60!black, linkcolor=blue!60!black]{hyperref}
\setlength{\parindent}{2em}
\setlength{\parskip}{1.25em}
\renewcommand{\baselinestretch}{1.0}
\begin{document}
\begin{center}
{\LARGE \textit{Tennessee v. Lane}, 541 U.S. 509 (2004)\par}
\vspace{0.8em}
{\large Case Brief --- Four-Page Edition\par}
\end{center}
\vspace{1em}

\section{Memory Jogger}

Congress validly exercised its Section~5 enforcement power to abrogate state sovereign immunity under Title~II of the Americans with Disabilities Act insofar as Title~II applies to the class of cases implicating the fundamental right of access to the courts (\href{https://www.law.cornell.edu/supremecourt/text/541/509}{541 U.S. 509}).

\section{Detailed Case Facts}

Title~II of the ADA provides that no qualified individual with a disability shall be excluded from or denied the benefits of a public entity's services, programs, or activities, and incorporates the Rehabilitation Act's damages remedies. Congress expressly provided that a State ``shall not be immune under the eleventh amendment'' from suit.

George Lane and Beverly Jones, paraplegics who use wheelchairs, sued Tennessee and several counties in 1998. Lane, summoned to answer criminal charges in a second-floor Polk County courtroom with no elevator, crawled up two flights of stairs for his first appearance. When he returned and refused to crawl or be carried, he was arrested and jailed for failure to appear. (The dissent noted that the court held one hearing in a first-floor library and offered to move proceedings to an accessible courthouse.) Jones, a certified court reporter, alleged that inaccessible courthouses cost her work and the chance to participate in the judicial process. Both sought damages and equitable relief.

\section{Procedural History}

The district court denied Tennessee's motion to dismiss on sovereign-immunity grounds, and Tennessee appealed; the United States intervened. The Sixth Circuit awaited \textit{Board of Trustees v.\ Garrett}, which barred Title~I damages suits against States but reserved Title~II, and then, under its en banc \textit{Popovich} decision allowing due-process-based Title~II claims, affirmed. The Supreme Court affirmed on May 17, 2004.

\section{Judicial Votes}

Affirmed, \textbf{5--4}.
\begin{itemize}
\item \textbf{Majority:} Justice Stevens, for Justices O'Connor, Souter, Ginsburg, and Breyer.
\item \textbf{Concurring:} Justice Souter, joined by Justice Ginsburg; Justice Ginsburg, joined by Justices Souter and Breyer.
\item \textbf{Dissenting:} Chief Justice Rehnquist, joined by Justices Kennedy and Thomas; Justice Scalia; Justice Thomas.
\end{itemize}

\section{Holding}

Under Section~5 and the congruence-and-proportionality test of \textit{City of Boerne v.\ Flores}, Title~II validly abrogates state immunity as applied to cases implicating the fundamental right of access to the courts. Courts need not assess a broad Section~5 statute as an undifferentiated whole but may examine it as applied to the category of rights at stake. The Court declined to decide Title~II's validity in other settings, such as hockey rinks, voting, education, or claims reviewed only under \textit{Cleburne}'s rational basis, so its broader survey of disability discrimination is context, not holding.

\section{Analysis of Opinions}

\textbf{Majority (Stevens).} Applying \textit{Kimel}, the Court found an unequivocal intent to abrogate and asked whether Congress acted under valid authority. Section~5 permits prophylactic legislation if it is congruent and proportional to the injury. Unlike Title~I in \textit{Garrett}, Title~II enforces fundamental rights subject to heightened scrutiny: access to the courts under due process, the Sixth Amendment's confrontation and jury-trial rights, and the public's First Amendment right of access. As in \textit{Hibbs}, heightened scrutiny made a pattern of violations easier to show. The record included decisions documenting unconstitutional treatment, a 1983 Civil Rights Commission finding that 76\% of public services in state buildings were inaccessible, and testimony on inaccessible courthouses. The remedy was proportional because it requires only reasonable modifications without undue burden or fundamental alteration, and older facilities may comply by relocating proceedings or assigning aides.

\textbf{Souter (concurring).} A broader history, including \textit{Buck v.\ Bell}, shows the judiciary itself endorsed eugenic exclusions.

\textbf{Ginsburg (concurring).} Citing \textit{Olmstead}, she stressed that equal citizenship requires accommodation, and rejected requiring Congress to indict each State separately.

\textbf{Rehnquist (dissenting).} The record showed little evidence that States, as opposed to local governments, denied court access; physical barriers alone do not violate due process when assistance is offered; and the as-applied method rigs the test for a statute that reaches everything from courthouses to stadiums.

\textbf{Scalia (dissenting).} He repudiated congruence and proportionality as a flabby test and read ``enforce'' to permit remedies only for actual violations, retaining broader review only for racial discrimination. \textbf{Thomas (dissenting)} disavowed reliance on \textit{Hibbs}.

\section{Examples: Future Applications}

\begin{enumerate}
\item \textbf{Same-Side (Valid): An Interpreter for a Parental-Rights Hearing.} A deaf parent is denied an interpreter at a termination hearing. Court access and the interest recognized in \textit{M.L.B.\ v.\ S.L.J.} fall squarely within \textit{Lane}.
\item \textbf{Same-Side (Valid): An Inaccessible Prison Cell.} A paraplegic inmate cannot use his toilet. Where Title~II remedies actual constitutional violations, abrogation is valid, as \textit{United States v.\ Georgia} held.
\item \textbf{Opposite-Side (Invalid): Stadium Seating.} A state-owned hockey rink lacks accessible seating. No fundamental right is involved, so \textit{Garrett} controls.
\item \textbf{Opposite-Side (Invalid): Occupational Licensing.} A blind applicant is denied a barbering license without accommodation. Licensing receives rational-basis review, and lower courts generally deny abrogation.
\item \textbf{Fence-Sitter: Extended Exam Time.} A state university refuses extra time to a student with a learning disability. Education is not a fundamental right under \textit{Rodriguez}, but \textit{Lane} noted a long history of discrimination in education, and courts are divided.
\end{enumerate}

\section{Critique}

\textbf{Originalist and textualist views.} Echoing Scalia, textualists argue that the power ``to enforce'' does not authorize prohibiting constitutional state conduct, and that congruence and proportionality is an atextual, indeterminate balancing test that invites judicial policymaking.

\textbf{Progressive and doctrinal views.} Disability-rights scholars welcomed the result and Ginsburg's account of accommodation as equal citizenship, but faulted the narrow, piecemeal holding. By tying validity to court access rather than revisiting \textit{Garrett} or treating disability as a suspect class, \textit{Lane} invited lower courts to deny abrogation in public services, health care, and education, forcing category-by-category litigation. The case pits state sovereign dignity and judicial supremacy against the Reconstruction Amendments' promise of national civil-rights enforcement.

\section{Key Quotations}

\begin{quote}
``For these reasons, we conclude that Title II, as it applies to the class of cases implicating the fundamental right of access to the courts, constitutes a valid exercise of Congress' \S~5 authority to enforce the guarantees of the Fourteenth Amendment.''
\end{quote}
\noindent --- Justice Stevens, Opinion of the Court, 541 U.S. at 533--534.

\begin{quote}
``Including individuals with disabilities among people who count in composing `We the People,' Congress understood in shaping the ADA, would sometimes require not blindfolded equality, but responsiveness to difference; not indifference, but accommodation.''
\end{quote}
\noindent --- Justice Ginsburg, concurring, 541 U.S. at 536.

\begin{quote}
``In conducting its as-applied analysis, however, the majority posits a hypothetical statute, never enacted by Congress, that applies only to courthouses. The effect is to rig the congruence-and-proportionality test by artificially constricting the scope of the statute to closely mirror a recognized constitutional right.''
\end{quote}
\noindent --- Chief Justice Rehnquist, dissenting, 541 U.S. at 551.

\end{document}
